\documentclass{article}

\textwidth=6.5in
\textheight=8.5in
\voffset=-54pt
\oddsidemargin=0in
\evensidemargin=0in
\setlength{\parskip}{8pt}
\setlength{\parindent}{0pt}

\usepackage{workbook}

\usepackage[sols]{handout}

\begin{document}

\htitle{Math Mb Final Review Session - Limits and Differential Equations}

{\large\textbf{Limits}}
  
\begin{enumerate}
 
    \item Evaluate the following limits or show why they do not exist.
    \begin{enumerate}
        \item $\displaystyle\lim_{x\to 1} \dfrac{\ln x}{e^x}$
        
        \begin{solution}
            We can use direct substitution:
            \[\displaystyle\lim_{x\to 1} \dfrac{\ln x}{e^x}=\dfrac{\ln 1}{e^1}=\boxed{0}\]
        \end{solution}
        \probonly{\vfill}
        
        \item $\displaystyle\lim_{x\to 2} \dfrac{4x-8}{x^2-4}$
        
        \begin{solution}
            As $x\to 2$, both the numerator and denominator go to $0$, so we should try to rewrite the expression by factoring.
            \begin{align*}
                \displaystyle\lim_{x\to 2} \dfrac{4x-8}{x^2-4} &= \displaystyle\lim_{x\to 2} \dfrac{4(x-2)}{(x+2)(x-2)}\\
                &= \displaystyle\lim_{x\to 2} \dfrac{4}{(x+2)}\\
                &= \boxed{1}
            \end{align*}
        \end{solution}
        \probonly{\vfill}
        
        \item $\displaystyle\lim_{x\to 1} \dfrac{e^x}{\ln x}$

        \begin{solution}
            As $x\to 1$, $e^x \to e$ and $\ln x \to 0$. To determine whether this limit is DNE, $\infty$ or $-\infty$, we need to look at the sign of the denominator. 
            
            As $x\to 1^-$, $\ln x$ is negative and approaching $0$, so $\displaystyle\lim_{x\to 1^-}\dfrac{e^x}{\ln x}=-\infty$. 

            As $x\to 1^+$, $\ln x$ is positive and approaching $0$, so $\displaystyle\lim_{x\to 1^+}\dfrac{e^x}{\ln x}=\infty$. 

            Since the two one-sided limits aren't equal, $\lim_{x\to 1}\dfrac{e^x}{\ln x}$ \textbf{does not exist}.
        \end{solution}
        \probonly{\vfill\vfill}
        
    \end{enumerate}

\probonly{\newpage}

  \begin{framed}
    \textbf{L'Hôpital's Rule.} \\ 
    If $\displaystyle \lim_{x \to a}
    \frac{f(x)}{g(x)}$ is a
    \cfitb[\vphantom{$\dfrac{0}{0}$}]{0.3in}{$\bm{{\dfrac{0}{0}}}$} or
    \cfitb[\vphantom{$\dfrac{0}{0}$}]{0.3in}{$\bm{{\dfrac{\infty}{\infty}}}$}
    type of limit, then $\displaystyle \lim_{x
      \to a} \frac{f(x)}{g(x)} =$
    \fitb[\vphantom{$\dfrac{0}{0}$}]{1.25in}{$\bm{\displaystyle \lim_{x \to a}
      \frac{f'(x)}{g'(x)}}$}, if the latter limit exists or is $\pm
    \infty$.\\

    Limits of the form ``$0\cdot \infty$'' and ``$1^\infty$'' can be rewritten to apply L'Hôpital's Rule.
  \end{framed}

\item Evaluate the following limits or show why they do not exist.
\begin{enumerate}
    
    \item $\lim\limits_{x \to \infty} \dfrac{\ln(x)}{x}$
        
        \begin{solution}
        This is an $\frac{\infty}{\infty}$ limit, so we can use
        L'H\^opital's Rule: 
        \[\displaystyle \lim_{x\to\infty} \frac{\ln
          x}{x} \lheq \lim_{x\to\infty} \frac{\quad \frac{1}{x} \quad}{1}
        = \boxed{0}\]
      \end{solution}
      \probonly{\vfill}
      
      \item $\lim\limits_{x \to \frac{\pi}{4}} \dfrac{\tan(x)-1}{4x - \pi}$
        
      \begin{solution}
        This is a $ \frac{0}{0}$ limit, so we can use L'H\^opital's Rule:
        \[\displaystyle \lim_{x\to\frac{\pi}{4}}\frac{\tan x - 1}{4x - \pi}
        \lheq \lim_{x\to\frac{\pi}{4}} \frac{\sec^2 x}{4} = \frac{2}{4} =
        \boxed{\frac{1}{2}}\]
      \end{solution}
      \probonly{\vfill}
        
        \item $\lim\limits_{x \to 0} \dfrac{e^{4x}-x-1}{e^x - 1}$
        
        \begin{solution}
                This is a $\frac{0}{0}$ limit, so we use L'H\^opital's Rule:
                \[\lim\limits_{x \to 0} \dfrac{e^{4x}-x-1}{e^x - 1} \lheq \lim\limits_{x \to 0} \dfrac{4e^{4x}-1}{e^x} = \dfrac{4-1}{1} = \boxed{3}\]
        \end{solution}
        \probonly{\vfill}
    
    \item $\lim\limits_{x \to \infty} \dfrac{\cos(x)}{x}$
        
        \begin{solution}
                Cosine oscillates between -1 and 1, so as $x \to \infty$, $\cos(x)$ does not have a limit. However, we have in the numerator something that is always in the interval $[-1,1]$ being divided by an increasing large number, and so  $\lim\limits_{x \to \infty} \dfrac{\cos(x)}{x} = \boxed{0}$.
        \end{solution}
        \probonly{\vfill}

    \probonly{\newpage}

    \item $\displaystyle \lim_{x \to -\infty} x e^x$
    
    \begin{solution}
    As $x \to -\infty$, $e^x \to 0$, so this is a $0 \cdot \infty$
    limit, and we should rewrite it so that we can use L'Hôpital's
    Rule:
    \begin{align*}
      \lim_{x \to -\infty} x e^x & \nolheq \lim_{x \to -\infty}
      \frac{x}{e^{-x}} \\
      & \lheq \lim_{x \to -\infty} \frac{1}{-e^{-x}} \\
      & \nolheq \boxed{0}
    \end{align*}
  \end{solution}
  \probonly{\vfill}

  \item $\displaystyle\lim_{x\to \infty} \dfrac{x}{\sin x}$
  
  \begin{solution}
      As $x\to \infty$, $\sin x$ oscillates between $-1$ and $1$ while the numerator $\to \infty$. When $\sin x$ is negative, $\dfrac{x}{\sin x}$ is a very large negative number, but when $\sin x$ is positive, $\dfrac{x}{\sin x}$ is a very large positive number, so the overall limit doesn't exist. 
  \end{solution}
  \probonly{\vfill}

  \item $\displaystyle\lim_{x\to\infty} x\sin\left(\dfrac{\pi}{x}\right)$
  
  \begin{solution}
      As $x\to\infty$, $\sin\left(\dfrac{\pi}{x}\right) \to \sin(0)=0$, so this is a ``$\infty \cdot 0$'' type of limit. Let's rewrite it so we can use L'Hôpital's Rule.
      \begin{align*}
          \displaystyle\lim_{x\to\infty} x\sin\left(\dfrac{\pi}{x}\right) &=\displaystyle\lim_{x\to\infty} \dfrac{\sin\left(\dfrac{\pi}{x}\right)}{1/x}\\
          &\lheq \displaystyle\lim_{x\to\infty} \dfrac{\cos\left(\dfrac{\pi}{x}\right)\cdot \dfrac{-\pi}{x^2}}{\dfrac{-1}{x^2}}\\
          &= \displaystyle\lim_{x\to\infty} \pi\cos\left(\dfrac{\pi}{x}\right)\\
          &=\boxed{\pi}
      \end{align*}
  \end{solution}
  \probonly{\vfill}


  \item $\displaystyle\lim_{x\to0^+} (\cos x)^{1/x^2}\vphantom{\frac{1}{1}}$

      \begin{solution} 
        As $x \to 0^+$, $\cos x \to 1$ and $\frac{1}{x^2} \to \infty$, so
        this is a $1^\infty$ limit.  Therefore, we should rewrite it so
        that we can apply L'Hôpital's Rule:
        \begin{align} 
          (\cos x)^{1/x^2} &= e^{\raisebox{3pt}{$\ln \left[ (\cos
                x)^{1/x^2} \right]$}} \nonumber \\
          &= e^{\raisebox{4pt}{$\frac{1}{x^2} \ln (\cos
              x)$}}
        \end{align}
        Now, let's focus on what happens to the exponent
        in the last expression as $x \to 0^+$:
        \begin{align*}
          \lim_{x \to 0^+} \frac{1}{x^2} \ln
          (\cos x) & \nolheq \lim_{x \to 0^+} \frac{\ln (\cos x)}{x^2} \\
          \shortintertext{This is a $\frac{0}{0}$ limit:}
          & \lheq \lim_{x \to 0^+} \frac{\frac{1}{\cos x} \cdot (-\sin
            x)}{2x} \\
          \intertext{This is still a $\frac{0}{0}$ limit; let's
            rewrite it before we apply L'Hôpital's Rule again:}
          & \nolheq \lim_{x \to 0^+} \frac{- \tan x}{2x} \\
          & \lheq \lim_{x \to 0^+} \frac{- \sec^2 x}{2} \\
          &= - \frac{1}{2}
        \end{align*}
        Remember that this was only the limit of the exponent in
        \eqref{eq:limit-1-to-the-inf-1}, so the limit we want is
        $\boxed{e^{-1/2}}$.
      \end{solution}
      \probonly{\vfill\vfill}
      
\end{enumerate}

\probonly{\newpage}

\begin{framed}
    Suppose $f(x)$ and $g(x)$ are functions that are both positive when
    $x$ is large.  We say that $g(x)$ \uline{grows more quickly than}
    $f(x)$, written $g(x) \gg f(x)$, if $\displaystyle \lim_{x \to
      \infty} \frac{f(x)}{g(x)} = \cfitb{0.25in}{0}$ (or, equivalently,
    if $\displaystyle \lim_{x \to \infty} \frac{g(x)}{f(x)} =
    \cfitb{0.25in}{\infty}$).
  \end{framed}

\item Put the following 3 types of functions in order from slowest growing to fastest.
\begin{itemize}
    \item $x^p$ where $p$ is a positive constant
    \item $b^x$ where $b$ is a constant, $b>1$
    \item $\log_b x$ where $b$ is a constant, $b>1$
\end{itemize}
\begin{solution}
    \[\log_b x \ll x^p \ll b^x\]
\end{solution}
\probonly{\vfill}


\item As $x$ increases without bound, the functions
    $\ln x$, $x^n$, and $e^x$ all increase without bound, but they
    increase at different rates.  As you learned in class, we describe
    this using the notation $\ll$.

    Here are some common functions listed according to their relative
    rates of growth:
    \[ 2^x \ll 3^x \ll 10^x \ll x^x \] Insert into this list the functions
    $e^x$, $x^3$, $\ln x$, $x$, and $\sqrt{x}$.

    (Having a general idea of the relative growth rates of functions such
    as polynomials, exponentials, and logarithms is a handy skill in many
    fields.)  
    
    
    \begin{solution}
    \[\ln x \ll \sqrt{x} \ll x \ll x^3 \ll 2^x \ll e^x \ll 3^x \ll 10^x \ll
    x^x\] 
  \end{solution}
  \probonly{\vfill}

  \item Evaluate the following limits.

   \begin{enumerate}
        \item $\displaystyle{\lim_{x\to \infty} \frac{3^x + x^{700} + x\ln x}{e^x + \sqrt{x}\ln x}}$
        
        \begin{solution}
          \begin{align*}
              \lim_{x\to \infty} \frac{3^x + x^{700} + x\ln x}{e^x + \sqrt{x}\ln x} &= \lim_{x\to \infty} \frac{3^x}{e^x}\\
              &= \lim_{x\to \infty} \left(\frac{3}{e}\right)^x\\
              &= \boxed{\infty}
          \end{align*}
        \end{solution}
        \probonly{\vfill}
        
        \item $\displaystyle{\lim_{x\to -\infty} \frac{3^x + x^{700} + \sin x}{e^x + x^{101} + \sqrt{|x|}}}$
        
        \begin{solution}
          \begin{align*}
          \lim_{x\to -\infty} \frac{3^x + x^{700} + \sin x}{e^x + x^{101} + \sqrt{|x|}} &= \lim_{x \to -\infty}\frac{x^{700}}{x^{101}}\\
          &= \lim_{x \to -\infty} x^{599} \\
          &= \boxed{-\infty}
          \end{align*}
        \end{solution}
        \probonly{\vfill}


    \end{enumerate}
\end{enumerate}

\newpage

{\large\textbf{Differential Equations}}

% Diff Eq problem's from Michael's 2024 review session

\begin{enumerate}
   

    \item Write the general solution to $\frac{dy}{dt} = ky$.   \vfill 
    
    \item Write the general solution to $\frac{dy}{dt} = t^2$. \vfill


\item \begin{problem}[\vfill\vfill]
    \$1000 is deposited in a bank account.  The account has a nominal
    annual interest rate of 6\%, compounded continuously.  Write a differential equation to model this situation. 
    
    Then, solve your differential equation to write a formula for $M(t)$. How much money is in the account after 3 years? 
  \end{problem}
  \begin{solution}
    Let's have $M = M(t)$ be the amount of money in the account at time
    $t$, where $t$ is measured in years and $t = 0$ represents the time of
    the initial deposit.  Then, $\frac{dM}{dt}$ represents the rate of
    change of money (with respect to time).

    Since in interest rate is 6\%, the instantaneous rate of change is 6\% of what is present. So, $\boxed{\frac{dM}{dt} = 0.06M}$.
     
    
  \end{solution}
  
\item  \label{moneybags-compound-continuously-diffeq}



    
  \begin{problem}[\vfill]
    \$1000 is deposited in a different bank account.  The account also has a nominal annual interest rate of 6\%, compounded continuously, but now money is being deposited continuously at a rate of \$600 per year. Write a differential equation to model the situation
  \end{problem}

  \begin{solution}
    We have now added a new source of change to our problem--money is ``flowing'' into the account at a constant rate of 600 dollars per year. Because the two sources of income are separate, we can simply add them together: 
    $\boxed{\frac{dM}{dt} = 0.06M + 600}$. 
  \end{solution}
%\end{enumerate} 

\item  \label{moneybags-continuous-compounding-deposit-solution} 


  
  \begin{problem}[\vfill\newpage]
    Which of the following is the solution of $\displaystyle \frac{dM}{dt}
    = 0.06M + 600$ satisfying $M(0) = 1000$?
    \begin{enumerate}[label=\protect\maybecircle{\roman*.},ref=\roman*]
    \plainitem $M(t) = 1000 e^{0.06 t} + 600t$
    \circleitem \label{moneybags-continuous-compounding-deposit-solution-right-choice}
      $M(t) = 11,000 e^{0.06 t} - 10,000$ 
    \plainitem $M(t) = (1000 + 600 t) e^{0.06 t}$
    \end{enumerate}
  \end{problem}

  \begin{solution}
    To answer this, we can simply check each choice and see whether it
    satisfies $\frac{dM}{dt} = 0.06 M + 600$ and $M(0) = 1000$. We do this by first taking each choice, and evaluating the left-hand side (LHS) and the right-hand side (RHS) of the differential equation separately to see if they are equal. (You can think of this as ``plugging" each choice of M(t) into the differential equation $[M(t)]'=0.06(M(t))+600$ and checking if it makes the equation work.)

    \begin{enumerate}[label=\roman*.]
    \item If $M(t) = 1000 e^{0.06 t} + 600 t$, then
      \begin{align*}
        \text{LHS: } &\frac{dM}{dt} = 0.06 \left( 1000 e^{0.06 t} \right) + 600 \\
        \text{RHS: }&0.06 M + 600 = 0.06 \left( 1000 e^{0.06 t} + 600 t \right) + 600
      \end{align*}
      Therefore, $\frac{dM}{dt} \neq 0.06 M + 600$, so this is not the
      correct choice.

    \item If $M(t) = 11,000 e^{0.06 t} - 10,000$, then
      \begin{align*}
        \text{LHS: }&\frac{dM}{dt} = 0.06 \cdot 11,000 e^{0.06 t} \\
        \text{RHS: }&0.06 M + 600  = 0.06 \left( 11,000 e^{0.06 t} - 10,000 \right) +
        600 = 0.06 \cdot 11,000 e^{0.06 t}
      \end{align*}
      So, in this case, $\frac{dM}{dt} = 0.06 M + 600$.  Moreover, $M(0) =
      11,000 - 10,000 = 1,000$.  So, \fbox{this is the correct choice}.

    \item If $M(t) = (1000 + 600 t) e^{0.06 t}$, then
      \begin{align*}
        \text{LHS: }&\frac{dM}{dt} = 0.06 (1000 + 600 t) e^{0.06 t} + 600 e^{0.06 t} \\
        \text{RHS: }&0.06 M + 600 = 0.06 (1000 + 600 t) e^{0.06 t} + 600
      \end{align*}
      These are not equal, so this is not the correct choice.

     
    \end{enumerate}
  \end{solution}
  

    \item

  
    
    \begin{problem}
    Which of the following is a solution to $\frac{dy}{dt} = 1 - y$?


    \begin{enumerate}[label=\protect\maybecircle{\roman*.}]
    \plainitem $y = Ce^{-t}$
    \plainitem $y = Ce^t -t$
    \plainitem $y = Ce^{-t} -1$
    \circleitem $y = Ce^{-t} +1$
    \plainitem $y = Ce^t +1$
  \end{enumerate}


  \end{problem}

  \begin{solution}
 $$ \begin{array}{|l|c|l|c|} \hline &&& \cr
 &\mbox{LHS: } \frac{dy}{dt} ~~&\mbox{  RHS: }  1-y & \cr \hline &&& \cr i. & -Ce^{-t} & 1-Ce^{-t} & \mbox{Not equal}\cr \hline &&& \cr ii. & Ce^t -1 & 1-(Ce^t -t) & \mbox{Not equal}\cr\hline &&& \cr iii. & -Ce^{-t}  & 1-(Ce^{-t} -1) = 2-Ce^{-t} & \mbox{Not equal}\cr\hline &&& \cr iv. &- Ce^{-t}  & 1-(Ce^{-t} +1) = -Ce^{-t}& \mbox{\bf{Equal!} }\cr \hline &&& \cr v. & Ce^t  & 1-(Ce^t +1) = -Ce^t & \mbox{Not equal}\cr \hline 
  \end{array}$$
 
  \end{solution}







\item \label{IV drug}
\begin{enumerate}
    \item{} 

\begin{problem}[\vfill]  A hospital patient is hooked up to an IV, which delivers medication into her body at a steady rate of 8 mg per hour. Her body processes the medication and breaks it down at a rate that is proportional to the amount present in her system. 

  If $A(t)$ is the amount of medication in her body at time $t$, write a differential equation for $A$ that reflects this information. 
  
\end{problem}
    \begin{solution}
      We want to write a differential equation about $\frac{dA}{dt}$, the
      instantaneous rate of change (in mg/hour) of medication in the patient's body.
      There are two factors that cause the amount of medication in her
      body to change: the fact that she's taking medication and the fact that
      her body is metabolizing it.  We can put these together as
      \[ 
      \frac{dA}{dt} = \left(
        \begin{minipage}{2in}
          \centering rate of increase of medication due to the patient
          taking it
        \end{minipage} \right) - \left(
        \begin{minipage}{2in}
          \centering rate of decrease of medication due to metabolism
        \end{minipage} \right)\] 
      We are told that the rate at which she metabolizes medication is
      proportional to $A$, the amount in her body, so
      \[ \text{rate at which the patient metabolizes medication} = k A \] for some
      constant $k$.  This means that  $\boxed{
        \frac{dA}{dt} = 8 - k A}$.

    \end{solution}

 \item
    \begin{problem}[\vfill]
      When the patient has 20 mg of medication in her body, the amount $A(t)$ is increasing at 6 mg per hour. Incorporate this information into your differential equation. 
    \end{problem}
\begin{solution}
When $A = 20, \frac{dA}{dt} = 6$. Plugging this into our differential equation gives $$6 = \frac{dA}{dt} = 8 - k A = 8-20k$$
Solving for $k$ gives $20k = 2$, or $k = 0.1$

Answer: $\boxed{
        \frac{dA}{dt} = 8 - 0.1 A}$
\end{solution}
    
  \item
    \begin{problem}[\vfill]
      One of the following is a possible formula for $A(t)$.  Which one?
      \begin{multicols}{4}
        \begin{enumerate}[label=\protect\maybecircle{\Alph*.}]
        \plainitem $(80 t + 8) e^{-0.1t}$
        \plainitem $20 e^{-0.1t}$
        \circleitem $80 - 20 e^{-0.1t}$
        \plainitem $8 - 20 e^{0.1t}$
        \end{enumerate}        
      \end{multicols}
    \end{problem}
 \begin{solution}
  We are looking for the choice that satisfies our differential
      equation $\frac{dA}{dt} = 8 - 0.1A$. In this problem, we are not given an initial condition. 
   

    \begin{enumerate}[label=\roman*.]
    \item If $A(t) = (80 t + 8) e^{-0.1t}$, then
      \begin{align*}
        \text{LHS: } &\frac{dA}{dt} = 80e^{-0.1t} -0.1 (80 t + 8) e^{-0.1t} = (80-8t-0.8)e^{-0.1t} \\
        \text{RHS: }& 8 - 0.1A  = 8 - 0.1(80 t + 8) e^{-0.1t}
      \end{align*}
      These are clearly not equal, so this is not the
      correct choice.

    \item If $A(t) = 20 e^{-0.1t}$, then
      \begin{align*}
        \text{LHS: }&\frac{dA}{dt} = -2 e^{-0.1t} \\
        \text{RHS: }& 8 - 0.1A = 8 -2 e^{-0.1t}
      \end{align*}
     These are clearly not equal, so this is not the
      correct choice.
      
    \item If $A(t) = 80 - 20 e^{-0.1t}$, then
      \begin{align*}
        \text{LHS: }&\frac{dA}{dt} = 2e^{-0.1t} \\
        \text{RHS: }&8 - 0.1A = 8  - 0.1\left(80 - 20 e^{-0.1t}\right) = 8 - 8 + 0.2e^{-0.1t} = 0.2e^{-0.1t}
      \end{align*}
      These are  equal, so \fbox{this is  the correct choice!}
\item (For completeness, we'll check the last one:) If $A(t) = 8 - 20 e^{0.1t}$, then
      \begin{align*}
        \text{LHS: }&\frac{dA}{dt} = -2 e^{0.1t} \\
        \text{RHS: }& 8 - 0.1A = 8  - 0.1\left(8 - 20 e^{0.1t}\right) = 8 - 0.8 + 0.2e^{0.1t}
      \end{align*}
     These are clearly not equal, so this is not the
      correct choice.
      
     
    \end{enumerate}
  \end{solution}
  
 
 
\end{enumerate}


\item{}
    The following model of blood cholesterol levels is based on the assumption
    that cholesterol is manufactured by the body for use in the
    construction of cell walls and is absorbed from foods containing
    cholesterol.  Let $C(t)$ be the amount (in mg per deciliter) of
    cholesterol in the blood of a particular person at time $t$ (in days).
    Because of the person's diet, their blood cholesterol satisfies the
    following differential equation:
    \[ \frac{dC}{dt} = 0.1 (200 - C) + 10 \] 

    \begin{enumerate}
        \item \begin{problem}[\vfill] Suppose that a person's cholesterol level is currently $C = 260$ mg per deciliter. Is their cholesterol going up or down? \end{problem}
        
      \begin{solution}
    ``Is their cholesterol going up or down?'' is really just a question about whether the derivative is positive or negative. 

    If $C = 260$, $\frac{dC}{dt} = 0.1 (200 - 260) + 10 = -6 + 10 = 4$, which is positive. So, the person's cholesterol is going {\bf up} at this moment. 
      \end{solution}
        
        
       

     \item \begin{problem}[\vfill] Still assuming that $C = 260$ today. Estimate the value of $C$ two days from now. 
\end{problem}
     
  \begin{solution}
   We calculated that $\frac{dC}{dt} = 4$--so the instantaneous rate of change is 4 mg per deciliter per day. 

   Thus, in two days we estimate that the cholesterol increases by $2 \times 4 = 8$. 

 \fbox{  After two days, the person's cholesterol is approximately 268 mg per deciliter. }
      \end{solution}
  
    \end{enumerate}

 \end{enumerate}


\end{document}